Atualmente, as aplicações têm se concentrado na coleta de informações a respeito dos usuários, com o objetivo de conhecer suas preferências, possibilitando a oferta de serviços personalizados. Assim, estas aplicações organizam seu conhecimento a respeito do usuário em perfis de usuário, os quais armazenam as informações em um formato baseado em um modelo de usuário específico.\cite{viviani10}

Sistemas sensíveis a contexto são aplicações nas quais o contexto atual de um usuário é levado em consideração na disponibilização de serviços. Dey\cite{dey2001} define contexto como qualquer informação que pode ser utilizada para caracterizar a situação de uma entidade, onde uma entidade pode ser uma pessoa, local ou objeto considerado relevante à interação entre o usuário e a aplicação. Desta forma, o perfil torna-se parte essencial na descoberta de informações pertinentes aos usuários que possam ser utilizadas na construção de um contexto.

Os ambientes de computação ubíqua são parte de um cenário no qual o custo de dispositivos é baixo e a capacidade e disponibilidade são altos, sendo possível prover serviços e informações de modo amplo e contínuo\cite{weiser91}. Esta visão vem sendo potencializada pela diminuição do custo e aumento da capacidade de memória e processamento dos dispositivos computacionais, bem como pela cada vez maior disponibilidade de recursos de localização e conectividade. A computação ubíqua enfatiza também a ``tecnologia calma'', onde a sobrecarga de informação é substituída pela disponibilização de recursos de forma inteligente, tomando a atenção do usuário apenas nos momentos em que é mais relevante.\cite{weiser96}

% TODO - usar a justificativa de tecnologia calma como justificativa para o trabalho

A capacidade de uma aplicação inferir o contexto de uma entidade pode ser amplificada através do uso de trilhas\cite{silva09}. Uma trilha é o histórico de contextos visitados por uma entidade, e das ações realizadas pela entidade nos contextos.

Existem trabalhos que utilizam trilhas como forma de armazenar e processar informações de um usuário\cite{silva08,silva09}. No entanto, nestes trabalhos, o modelo do usuário já está pré-definido, e a trilha é utilizada apenas para alimentar informações específicas do perfil.

Este trabalho busca aprofundar esta abordagem apresentando um modelo para gerenciamento de perfis através de inferência em trilhas. O objetivo do modelo é permitir que qualquer aplicação que faça uso de trilhas possa ser capaz de montar seu próprio modelo de usuário e definir as regras necessárias para que este perfil seja inferido a partir das trilhas.
